\scanpage{68}
% letterhead
Séminaire Nicolas BOURBAKI

19, 20, 21 juin 1976 - Amphi. Hermite - 14 H 30

Institut Henri Poincaré - 11 rue P. et M. Curie

\uline{PARIS 5e}

\uline{Samedi 19 juin}

A. GRAMAIN - Rapport sur la théorie des noeuds.

A. GUICHARDET - Représentations des groupes $\mathrm G^{\mathrm X}$ d'après Gelfand et Delorme.

\uline{Dimanche 20 juin}

D. FERRAND - Les modules projectifs sur un anneau de polynômes sur un corps sont libres, d'après Quillen et Suslin.

J.-P. SERRE - Représentations linéaires des groupes finis \textquotedbl{}algébriques\textquotedbl{} (d'après Deligne-Lusztig).

\uline{Lundi 21 juin} (Première Séance à \uline{14 Heures})

M. WALDSCHMIDT - Les travaux de C.V. \handwritten{\rlap{\v{\phantom{C}}}}Cudnovskij sur les nombres transcendants.

M. DEMAZURE - Identités de Macdonald.

Dates des prochaines réunions du séminaire :

20-22 novembre 1976,

19-21 février 1977,

18-20 juin 1977.

% archive sheet number: 67

